\begin{abstract}
    \chapter{Abstract}

        Multi-agent systems by their nature have synchronization and communication limitations that consequently render them challenging to control. A reliable solution is the problem of optimal consensus. Each agent is assigned with a cost function of the output vector, which is transmitted to neighboring agents through the topology of the communication graph. The goal is to drive the outputs of the agents to a common value that minimizes the sum of the cost functions. In this thesis, are assumed non-convex cost functions with connected topologies for the communication graph. The transmission of the position vectors between agents is done at specific sampling periods, in order to secure the information from system disturbances. A regulating algorithm was designed for this task through the NOCGPROX variables, which are the agent's position errors. Taking into consideration the dynamics of every agent, distributed robust control laws were designed to ensure that the outputs approach the optimal point or a region of it. These laws were developed with the sliding surface method to ensure robustness and feedback linearization to cancel the nonlinear dynamics of the system. These laws ensure that these variables tend to zero after a number of iterations, along with its square integration, for the precise implementation of this algorithm. The above algorithm and laws were applied to a swarm of UAVs. A number of simulation scenarios were conducted to ensure the accurate implementation and execution of the algorithm and control laws in MATLAB.

        \vspace{20mm}
        \leftline{\large\textbf{Keywords}}
        \normalsize
        \vspace{5mm}
        \begin{flushleft}
            Optimal Consensus, Distributed Non-Convex Optimization, Multi-Agent Systems, Asynchronous Communication, Distributed Automatic Control.
        \end{flushleft}
\end{abstract}


\begin{abstract}
    \chapter{Περίληψη}
         Τα πολυπρακτορικά συστήματα από τη φύση τους έχουν προβλήματα συγχρονισμού και επικοινωνιών που καθιστούν δύσκολο τον έλεγχό τους. Μία αξιόπιστη λύση αποτελεί το πρόβλημα της βέλτιστης συναίνεσης. Σε κάθε πράκτορα αντιστοιχεί μια συνάρτηση κόστους του διανύσματος εξόδου το οποίο μεταδίδεται στους γειτονικούς πράκτορες μέσω της τοπολογίας του γράφου επικοινωνίας. Ο στόχος είναι οι έξοδοι των πρακτόρων να οδηγηθούν σε κοινή τιμή που ελαχιστοποιεί το άθροισμα των συναρτήσεων κόστους. Στην εργασία αυτή υποθέτουμε μη κυρτές συναρτήσεις κόστους με συνδεδεμένη τοπολογία για τον γράφο επικοινωνίας. Η μετάδοση των θέσεων μεταξύ των πρακτόρων γίνεται σε συγκεκριμένες περιόδους δειγματοληψίας για θωράκιση της πληροφορίας από διαταραχές του συστήματος. Για τη διαδικασία αυτή σχεδιάστηκε ένας ρυθμιστικός αλγόριθμος, μέσω των μεταβλητών NOCGPROX, οι οποίοι αποτελούν τα σφάλματα θέσεις των πρακτόρων. Λαμβάνοντας υπόψη τη δυναμική των πρακτόρων, σχεδιάστηκαν κατανεμημένοι εύρωστοι νόμοι ελέγχου που εξασφαλίζουν ότι οι έξοδοι προσεγγίζουν το βέλτιστο σημείο ή μια περιοχή του. Οι νόμοι αυτοί βασίστηκαν στη μεθοδολογία επιφάνειας ολίσθησης για να υπάρχει ευρωστεία, σε συνδυσμό με γραμμικοποίηση με ανάδραση για να ακυρωθούν οι μη γραμμικές δυναμικές του συστήματος. Οι νόμοι αυτοί εξασφαλίζουν τον μηδενισμό των μεταβλητών αυτών, μαζί και την τετραγωνική του ολοκλήρωση, για την ορθή λειτουργία του αλγορίθμου. Ο παραπάνω αλγόριθμος και οι νόμοι εφαρμόστηκαν σε ένα σμήνος από μη επανδρωμένων αεροσκαφών. Διάφορα σενάρια προσομοιώσεων εκτελέστηκαν με στόχο την εξασφάλιση της ορθής λειτουργίας του αλγορίθμου και των νόμων ελέγχου, μέσω του περιβάλλοντος του MATLAB.
        
    \vspace{20mm}
    \leftline{\large\textbf{Λέξεις Κλειδιά}}
    \vspace{5mm}
    \begin{flushleft}
        Βέλτιστη Συναίνεση, Μη Κυρτή Κατανεμημένη Βελτιστοποίηση, Πολυπρακτορικά Συστήματα, Ασύγχρονες Επικοινωνίες, Κατανεμημένος Αυτόματος Έλεγχος.
     \end{flushleft}
\end{abstract}
